\section{Lifecycle and State Machines}

The parent execution lifecycle is monotonic. Terminal executions do not return to \status{RUNNING}. The engine serializes operations per execution so API requests, background ticks, stream events, manual cancel, and reconciliation cannot concurrently mutate the same execution record.

\begin{figure}[t]
\centering
\small
\begin{tabularx}{\linewidth}{>{\raggedright\arraybackslash}p{0.18\linewidth} >{\raggedright\arraybackslash}p{0.28\linewidth} X}
\toprule
Phase & Allowed next phase & Purpose \\
\midrule
\status{CREATED} & \status{VALIDATING} & Request accepted and domain objects built. \\
\status{VALIDATING} & \status{RUNNING} or terminal & Position, symbol rules, bounds, duration, and dust checks. \\
\status{RUNNING} & \status{CANCELLING} or terminal & Submit, cancel, reconcile, deadline, and fill processing. \\
\status{CANCELLING} & terminal & Stop active exposure, reconcile, then summarize. \\
Terminal & none & \status{COMPLETED}, \status{PARTIALLY COMPLETED}, \status{EXPIRED}, \status{CANCELLED}, or \status{FAILED}. \\
\bottomrule
\end{tabularx}
\caption{Parent execution state transition table. Terminal states are one-way.}
\label{fig:execution-state}
\end{figure}

\begin{figure}[t]
\centering
\begin{tikzpicture}[
  node distance=0.48cm and 0.6cm,
  state/.style={draw, rounded corners=2pt, align=center, minimum width=2.15cm, minimum height=0.55cm, font=\small},
  terminal/.style={state, fill=gray!10},
  arrow/.style={-{Latex[length=2mm]}, thick}
]
\node[state] (pending) {PENDING\\SUBMIT};
\node[state, right=of pending] (open) {OPEN};
\node[state, right=of open] (partial) {PARTIALLY\\FILLED};
\node[state, right=of partial] (pcancel) {PENDING\\CANCEL};
\node[terminal, below=0.65cm of pending] (unknown) {UNKNOWN};
\node[terminal, below=0.65cm of open] (rejected) {REJECTED};
\node[terminal, below=0.65cm of partial] (filled) {FILLED};
\node[terminal, below=0.65cm of pcancel] (cancelled) {CANCELLED};
\draw[arrow] (pending) -- (open);
\draw[arrow] (pending) -- (unknown);
\draw[arrow] (pending) -- (rejected);
\draw[arrow] (open) -- (partial);
\draw[arrow] (open) -- (pcancel);
\draw[arrow] (open) -- (filled);
\draw[arrow] (partial) -- (pcancel);
\draw[arrow] (partial) -- (filled);
\draw[arrow] (pcancel) -- (cancelled);
\draw[arrow] (pcancel) -- (filled);
\end{tikzpicture}
\caption{Child order states. \status{UNKNOWN} is a reserved-exposure state, not a harmless error.}
\label{fig:child-state}
\end{figure}

The child state machine is designed for disagreement between REST responses and user-stream events. A fill can arrive after a cancel request. A create request can time out while the exchange still accepts the order. A REST snapshot can be older than a private stream event. The engine therefore treats cumulative fills as monotonic, keeps pending-cancel and unknown-create exposure reserved, and uses exact client-order lookup to repair ambiguous create outcomes.

\textbf{Auditability.} Parent execution IDs, internal child order IDs, and Binance client order IDs are intentionally related. A parent execution such as \code{exec\_3168600ee25b4193} produces client order IDs with a short deterministic prefix such as \code{ce\_3168600ee25b\_1}. This makes logs, child-order CSVs, fill CSVs, reconciliation rows, and Binance order evidence joinable without relying on timestamps alone. It also scopes reconciliation to the execution under review; broad prefix scans can help diagnostics, but exact client-order lookup is the proof used for ambiguous create outcomes.
